\documentclass[11pt]{article}
\usepackage[margin=1in]{geometry}
\usepackage{amsmath,amssymb,amsthm,mathtools}
\usepackage{graphicx}
\usepackage{microtype}
\usepackage{enumitem}
\usepackage[colorlinks=true,linkcolor=blue,citecolor=blue,urlcolor=blue]{hyperref}

\newtheorem{theorem}{Theorem}
\newtheorem{lemma}{Lemma}
\newtheorem{remark}{Remark}
\newcommand{\1}{\mathbf 1}
\newcommand{\esssup}{\operatorname*{ess\,sup}}
\newcommand{\essinf}{\operatorname*{ess\,inf}}
\newcommand{\supp}{\operatorname{supp}}

\title{Every amenable action admits a one-approximating sequence\\
\large Candidate full solution to the conjecture after Definition 1.3 of arXiv:1611.08486}
\author{Autonomous proof packet --- agent\_lane\_09}
\date{13 August 2026}

\begin{document}
\maketitle

\begin{center}
\fbox{\parbox{0.92\textwidth}{
\textbf{Status.} Candidate full solution, likely valid, pending expert review.
The proof covers locally compact Hausdorff groups in the net convention used by
the source paper.  No unimodularity assumption is made.
}}
\end{center}

\section{Source question}

The source is A. M. Gonz\'alez-P\'erez, \emph{Crossed-Product Extensions of
$L_p$-Bounds for Amenable Actions}, arXiv:1611.08486, published in
\emph{J. Functional Analysis} 274 (2018), 2846--2883
\cite{GonzalezPerez}.  On page 9, after Definition 1.3, it asks whether every
amenable action has a $C$-approximating sequence with $C=1$.

\begin{figure}[ht]
  \centering
  \includegraphics[width=\textwidth]{figures/open_problem_crop.png}
  \caption{Definition 1.3 and the conjecture on source-paper page 9.}
\end{figure}

Write the centre of the semifinite von Neumann algebra as
$\mathcal Z(\mathcal M)=L_\infty(X)$.  For a measurable
$F\subset X\times G$, put
\[
  F_x=\{h\in G:(x,h)\in F\}.
\]
The case $C=1$ means that $\mu(F_x)$ is independent of $x$ (up to null sets),
and the normalized indicators
\[
  z_F(x,h)=\mu(F_x)^{-1/2}\1_{F_x}(h)
\]
must be asymptotically equivariant.  When all fibres have measure $L$, the
squared fibrewise $L_2$ defect is exactly
\begin{equation}\label{eq:defect-identity}
 \int_G\left|z_F(g^{-1}x,g^{-1}h)-z_F(x,h)\right|^2\,d\mu(h)
 =\frac{\mu\bigl(gF_{g^{-1}x}\triangle F_x\bigr)}{L}.
\end{equation}

\section{Proof intuition}

Amenability supplies positive compactly supported fibre functions $u_x$ of
$L_2$ norm one.  Their squares $p_x=u_x^2$ are probability densities with
small $L_1$ equivariance defect.  The subgraph of $p_x$ would be an
equal-volume set, but it lives in $G\times\mathbb R_+$ rather than in $G$.
We replace the vertical coordinate by finitely many \emph{right-translation
slots}: quantized level sets of $p_x$ are placed in mutually disjoint sets
$A_j(x)r_j\subset G$.  Left translation commutes with these right slots, so
the layer-cake estimate controls their symmetric-difference defect.

The finite layer sum misses its target volume by a bounded amount depending on
$x$.  One extra slot supplies exactly that missing measure.  This correction
can behave arbitrarily under the action: its measure is uniformly bounded,
whereas the total fibre measure tends to infinity, so its relative defect
tends to zero.

Right translation need not preserve left Haar measure.  This is overcome by
placing every slot in $H=\ker\Delta_G$.  If $G$ is nonamenable, then $H$ is
noncompact: otherwise $H$ would be amenable and the abelian quotient $G/H$
would make $G$ an amenable extension.  Thus $H$ supplies arbitrarily many
disjoint Haar-preserving slots.  If $G$ itself is amenable, ordinary F\o lner
sets already solve the problem for every action.

\section{Statement}

We use the source paper's characterization of amenability: there is a net of
positive $u_\alpha\in C_c(G;\mathcal Z(\mathcal M))$ such that
\begin{align}
 \int_G u_\alpha(x,h)^2\,d\mu(h)&=1,\label{eq:norm-one}\\
 \esssup_x\int_G
 \left|u_\alpha(g^{-1}x,g^{-1}h)-u_\alpha(x,h)\right|^2d\mu(h)&\longrightarrow0
 \quad(g\in G).\label{eq:l2-amen}
\end{align}
As in the source, ``sequence'' permits a net; under the usual countability
hypotheses one may diagonalize to a literal sequence.

\begin{theorem}\label{thm:main}
Let $G$ be a locally compact Hausdorff group with left Haar measure $\mu$, and
let $\theta:G\to\operatorname{Aut}(\mathcal M)$ be an amenable action in the
sense of \eqref{eq:norm-one}--\eqref{eq:l2-amen}.  Then there is a net of
measurable sets $F_\beta\subset X\times G$ such that, for every $\beta$,
\[
  1\leq \esssup_x\mu((F_\beta)_x)
   =\essinf_x\mu((F_\beta)_x)<\infty,
\]
and the normalized indicators satisfy \eqref{eq:l2-amen}.  Hence every
amenable action admits a $1$-approximating sequence in the sense of Definition
1.3 of \cite{GonzalezPerez}.
\end{theorem}

\section{Two elementary lemmas}

Let $\Delta_G:G\to(0,\infty)$ be the modular homomorphism and put
$H=\ker\Delta_G$.  Right translation by any element of $H$ preserves $\mu$.

\begin{lemma}[Haar-preserving slots]\label{lem:slots}
If $H$ is noncompact, then for every compact $D\subset G$ and every
$n\in\mathbb N$ there are $r_0,\ldots,r_n\in H$ such that the sets
$Dr_0,\ldots,Dr_n$ are pairwise disjoint.
\end{lemma}

\begin{proof}
Suppose $r_0,\ldots,r_{k-1}$ have been selected.  The intersection
\[
 H\cap\bigcup_{i<k}D^{-1}Dr_i
\]
is compact.  Choose $r_k$ outside it, which is possible because $H$ is
noncompact.  If $Dr_k$ met $Dr_i$, then
$r_k\in D^{-1}Dr_i$, a contradiction.
\end{proof}

\begin{lemma}[Equal-volume layer encoding]\label{lem:encoding}
Assume $H$ is noncompact.  Let $Q\subset G$ be finite and let
$p:X\times G\to[0,\infty)$ be measurable, supported in a common compact set
$S\subset G$, with
\[
 \int_Gp(x,h)\,d\mu(h)=1,\qquad 0\leq p(x,h)\leq M.
\]
For every sufficiently large $N$ there is a measurable
$F\subset X\times G$ whose fibres all have one finite measure $L\geq1$ and,
for every $g\in Q$,
\begin{equation}\label{eq:encoding-bound}
 \esssup_x\frac{\mu(gF_{g^{-1}x}\triangle F_x)}{L}
 \leq \esssup_x\|g p_{g^{-1}x}-p_x\|_1+C_{S}\frac{M}{N},
\end{equation}
where $C_S<\infty$ is independent of $N$, $x$, and $g\in Q$.
Here $(g p_{g^{-1}x})(h)=p(g^{-1}x,g^{-1}h)$.
\end{lemma}

\begin{proof}
Set $\delta=M/N$ and
\[
 A_j(x)=\{h:p(x,h)\geq j\delta\},\qquad 1\leq j\leq N,
 \qquad V(x)=\sum_{j=1}^N\mu(A_j(x)).
\]
The functions $x\mapsto\mu(A_j(x))$ are measurable.  Pointwise quantization
gives
\begin{equation}\label{eq:volume-deficit}
 0\leq 1-\delta V(x)<\delta\mu(S).
\end{equation}

First suppose $G$ is nondiscrete.  Haar measure is nonatomic.  Take
$L=\delta^{-1}$ and $d(x)=L-V(x)$.  By \eqref{eq:volume-deficit},
$0\leq d(x)<\mu(S)$.  In a fixed compact measurable set $R$ of measure at
least $\mu(S)$, choose a nested measurable family $(R_t)_{0\leq t\leq\mu(S)}$
with $\mu(R_t)=t$.  Such a family follows, for example, by recursive bisection
of the nonatomic finite measure space.  Put $C_x=R_{d(x)}$.

If $G$ is discrete, normalize Haar measure to counting measure.  Put
$L=\lceil\delta^{-1}\rceil$.  Since $V(x)$ is integral,
$d(x)=L-V(x)$ is an integer and \eqref{eq:volume-deficit} implies
\[
 0\leq d(x)\leq |S|+1.
\]
Choose a fixed finite ordered set $R$ with at least $|S|+1$ points and let
$C_x$ consist of its first $d(x)$ points.  This also makes $(x,h)\mapsto
\1_{C_x}(h)$ measurable.

Let $c_S=\mu(S)$ in the nondiscrete case and $c_S=|S|+1$ in the discrete
case.  Choose a compact set $D$ containing
\[
 \bigcup_{g\in Q\cup\{e\}}g(S\cup R).
\]
By Lemma \ref{lem:slots}, choose $r_0,\ldots,r_N\in H$ for which the
$Dr_j$ are disjoint, and define
\begin{equation}\label{eq:F-construction}
 F_x= C_xr_0\ \cup\ \bigcup_{j=1}^N A_j(x)r_j.
\end{equation}
All right translations in \eqref{eq:F-construction} preserve Haar measure,
the pieces are disjoint, and therefore
\[
 \mu(F_x)=\mu(C_x)+V(x)=d(x)+V(x)=L
\]
for every $x$.

Fix $g\in Q$.  The enlarged translated pieces occurring in
$gF_{g^{-1}x}$ and $F_x$ still lie in the mutually disjoint $Dr_j$ slots.
Consequently
\begin{align}
 \mu(gF_{g^{-1}x}\triangle F_x)
 &=\sum_{j=1}^N\mu(gA_j(g^{-1}x)\triangle A_j(x))
   +\mu(gC_{g^{-1}x}\triangle C_x).\label{eq:split}
\end{align}
For $a,b\in[0,M]$,
\begin{equation}\label{eq:scalar-quantization}
 \delta\sum_{j=1}^N
 \left|\1_{\{a\geq j\delta\}}-\1_{\{b\geq j\delta\}}\right|
 \leq |a-b|+\delta.
\end{equation}
Apply this at
$a=p(g^{-1}x,g^{-1}h)$ and $b=p(x,h)$ and integrate.  Since the integrand
vanishes off $gS\cup S$,
\begin{equation}\label{eq:main-layers}
 \delta\sum_{j=1}^N\mu(gA_j(g^{-1}x)\triangle A_j(x))
 \leq\|g p_{g^{-1}x}-p_x\|_1+\delta\mu(gS\cup S).
\end{equation}
The correction term in \eqref{eq:split} is at most $2c_S$.  In both cases
$L\geq\delta^{-1}$, hence $L^{-1}\leq\delta$.  Divide
\eqref{eq:split} by $L$ and use \eqref{eq:main-layers} to obtain
\[
 \frac{\mu(gF_{g^{-1}x}\triangle F_x)}{L}
 \leq\|g p_{g^{-1}x}-p_x\|_1
   +\delta\bigl(\mu(gS\cup S)+2c_S\bigr).
\]
Left invariance gives $\mu(gS\cup S)\leq2\mu(S)$, proving
\eqref{eq:encoding-bound} with
$C_S=2\mu(S)+2c_S$.
\end{proof}

\section{Proof of the theorem}

\begin{proof}[Proof of Theorem \ref{thm:main}]
We divide into two cases.

\smallskip
\noindent\textbf{Case 1: $G$ is amenable.}
The F\o lner--Reiter criterion gives, for every finite $Q\subset G$ and
$\varepsilon>0$, a relatively compact measurable set $A$ of positive finite
measure such that
\[
 \frac{\mu(gA\triangle A)}{\mu(A)}<\varepsilon\qquad(g\in Q).
\]
For completeness, this follows by choosing
$f\in C_c(G)_+$ with $\int f=1$ and
$\sum_{g\in Q}\|gf-f\|_1<\varepsilon$, applying layer cake to
$A_t=\{f>t\}$, and selecting one level at which the ratio of the summed
boundaries to $\mu(A_t)$ is small.  In a noncompact group, $A$ can be replaced
by a right translate or a disjoint union of right translates so that its
measure is at least one without changing these ratios.  Indeed, if
$\ker\Delta_G$ is noncompact, use finitely many disjoint right translates from
that kernel.  If $\ker\Delta_G$ is compact, then noncompactness of $G$ forces
the right-Haar scaling factors to be unbounded, and one right translate
suffices.  In a compact group take $A=G$ with the standard probability
normalization of Haar measure.  Put $F_x=A$ for every $x$.  Then
\eqref{eq:defect-identity} tends to zero, independently of the action, and all
fibres have the same measure.

\smallskip
\noindent\textbf{Case 2: $G$ is nonamenable.}
The quotient $G/H$ is abelian, hence amenable.  Amenability is
stable under extensions \cite[Chapter 2]{Paterson}; thus $H$ cannot be
amenable.  In particular $H$ is noncompact.

Fix a finite $Q\subset G$ and $\varepsilon>0$.  By amenability of the action,
choose one positive compactly supported approximant $u$ satisfying
\eqref{eq:norm-one} and
\begin{equation}\label{eq:u-small}
 \esssup_x\|g u_{g^{-1}x}-u_x\|_2<\varepsilon/4
 \qquad(g\in Q).
\end{equation}
Put $p_x=u_x^2$.  It has integral one, a common compact support $S$, and a
finite uniform bound $M$.  Cauchy--Schwarz and \eqref{eq:norm-one} give
\begin{align}
 \|g p_{g^{-1}x}-p_x\|_1
 &=\int_G |(g u_{g^{-1}x})^2-u_x^2|\,d\mu\notag\\
 &\leq\|g u_{g^{-1}x}-u_x\|_2
       \|g u_{g^{-1}x}+u_x\|_2
 \leq2\|g u_{g^{-1}x}-u_x\|_2.\label{eq:square-bound}
\end{align}
Choose $N$ so large that the second term on the right of
\eqref{eq:encoding-bound} is less than $\varepsilon/2$ and that the resulting
common fibre measure is at least one.  Lemma \ref{lem:encoding},
\eqref{eq:u-small}, and \eqref{eq:square-bound} now yield
\[
 \esssup_x\frac{\mu(gF_{g^{-1}x}\triangle F_x)}{L}<\varepsilon
 \qquad(g\in Q).
\]
By \eqref{eq:defect-identity}, the normalized indicators have squared
$L_2$ defect less than $\varepsilon$ on $Q$.  Direct the construction by the
pairs $(Q,\varepsilon)$, ordered by increasing $Q$ and decreasing
$\varepsilon$.  This produces the required $1$-approximating net.
\end{proof}

\section{Verification report}

The proof was audited through the following independent checkpoints.

\begin{enumerate}[leftmargin=*]
\item \textbf{Modular dichotomy.}  If $G$ is nonamenable, then
$\ker\Delta_G$ is nonamenable and hence noncompact, because the quotient is
abelian and amenability is extension-stable.
\item \textbf{Slot separation.}  At the $k$th choice only the compact set
$H\cap\bigcup_{i<k}D^{-1}Dr_i$ is forbidden.  This cannot exhaust noncompact
$H$.
\item \textbf{Exact volume.}  The level sum has deficit at most $\mu(S)$
(or at most $|S|+1$ in the discrete case).  The correction layer fills that
deficit exactly, and every slot preserves Haar measure.
\item \textbf{Defect estimate.}  Equation \eqref{eq:scalar-quantization}
is the floor-function bound
$\delta|\lfloor a/\delta\rfloor-\lfloor b/\delta\rfloor|
\leq |a-b|+\delta$.  The arbitrary correction costs at most $2c_S/L$.
\item \textbf{Measurability.}  Tonelli gives measurability of all level
volumes.  A nonatomic finite measure space admits a nested family of every
intermediate measure; composing that family with the measurable deficit gives
a measurable correction field.  The discrete selection is finite-valued.
\item \textbf{Finite regression check.}  Running
\texttt{python code/verify\_quantized\_layers.py} tests 50,000 random scalar
floor inequalities and 5,000 random discrete volume corrections.  It prints
\texttt{PASS}.  This check is not used as proof.
\end{enumerate}

\section{Scope, limitations, and novelty search}

The result settles the $C=1$ conjecture as stated in the source paper, using
the source paper's own approximating-net characterization of an amenable
action.  It does not address the distinct conjecture later in that paper about
whether an intertwining complete isometry characterizes amenability.

The numerical lower bound ``$1\leq\mu(F_x)$'' depends on the arbitrary scalar
normalization of Haar measure.  The construction makes $L$ arbitrarily large
in the nonamenable case; the standard F\o lner amplification does the same for
noncompact amenable groups, while compact Haar measure is conventionally
normalized to one.

A bounded novelty check was performed on 12--13 August 2026.  It searched the
run's registry, solution, attempt, and proof-gap indexes; exact arXiv id and
paper-title queries; and the phrases ``every amenable action admits a
1-approximating sequence'', ``$C$-approximating sequences'', and
``accretivity constant'' together with amenable action/crossed product.  The
search also inspected the citation neighborhood containing arXiv:2009.03019
and later multiplier-transference work.  No paper claiming this conjecture was
found.  A 2026 thesis by the source author still records examples with
$C=1$ rather than a general theorem.  This is not an exhaustive
MathSciNet/zbMATH citation review, so mathematical confidence is high but
novelty confidence is moderate pending expert bibliographic review.

\section{Human-review recommendation}

Review as a candidate full solution.  The most valuable checks are: (i) the
measurable nested correction family in the exact measure-theoretic framework
intended by Definition 1.3; (ii) the harmless net/sequence convention; and
(iii) the standard extension theorem used to force noncompactness of
$\ker\Delta_G$.  No computational or unproved mathematical dependency remains.

\begin{thebibliography}{9}
\bibitem{GonzalezPerez}
A. M. Gonz\'alez-P\'erez,
\emph{Crossed-Product Extensions of $L_p$-Bounds for Amenable Actions},
J. Funct. Anal. 274 (2018), no. 10, 2846--2883;
\href{https://arxiv.org/abs/1611.08486}{arXiv:1611.08486}.

\bibitem{Paterson}
A. L. T. Paterson,
\emph{Amenability}, Mathematical Surveys and Monographs 29,
American Mathematical Society, 1988.
\end{thebibliography}

\end{document}
